% Source: https://artofproblemsolving.com/community/c3h2239757_birthday_cakes
% Stable identity: aops-topic-2239757-q01
\textbf{(a)} The cakes are eaten at the times $0,\ \dfrac{1}{2},\ \dfrac{3}{4},\ \dfrac{7}{8},\ldots$ (in years), all before the next birthday, and the amounts are $1,\dfrac{1}{2},\dfrac{1}{4},\dfrac{1}{8},\ldots$. After the first $k+1$ cakes the total is
\[1+\dfrac{1}{2}+\cdots+\dfrac{1}{2^k}=2-\dfrac{1}{2^k},\]
which approaches $2$. So the total is $\boxed{2}$ cakes (infinitely many pieces, eaten closer and closer to the next birthday).

\textbf{(b)} The moments of level $k\ge1$ are the times $\dfrac{m}{2^k}$ with $m$ odd and $0<m<2^k$, and there are $2^{k-1}$ of them. At each one you eat $\dfrac{1}{2^k}$ of a cake, so level $k$ contributes
\[2^{k-1}\cdot\dfrac{1}{2^k}=\dfrac{1}{2}\]
of a cake. With levels $1,2,\ldots,n$ the total is
\[\boxed{1+\dfrac {n}{2}}\ \text{cakes}.\]
This grows without bound as $n$ grows: if the process never stops, you would eat infinitely much cake in a single year, although each single piece gets smaller and smaller. (In (a) there is only one piece at each step, so the halving wins and the total stays below $2$; in (b) the number of pieces doubles from each level to the next, exactly cancelling the halving of their size.)
